\documentclass[../../main.tex]{subfiles}
\begin{document}

\begin{center}
  {\fontsize{48pt}{58pt}\selectfont \textbf{1991}\normalsize\textbf{年度}}
\end{center}

\begin{table}[htb]
  \centering
  \caption{出題テーマと難易度}
  \begin{tabular}{|c|p{9cm}|c|}
    \hline
    \textbf{問題番号} & \textbf{テーマ} & \textbf{難易度} \\
    \hline
    第1問 & 三角関数で定まる曲線と面積 & \\
    \hline
    第2問 & 正八面体の高さと注水量 & \\
    \hline
    第3問 & 3本の単位ベクトルのなす角と等号条件 & \\
    \hline
    第4問 & 二次曲線を写す整数行列 & \\
    \hline
    第5問 & 重複なしの抽出と期待値 & \\
    \hline
    第6問 & 指数関数の接線による不等式 & \\
    \hline
  \end{tabular}
\end{table}

\begin{table}[htb]
  \centering
  \caption{解答}
  \renewcommand{\arraystretch}{1.4}
  \begin{tabular}{|c|c|p{9cm}|}
    \hline
    \textbf{問題番号} & \textbf{小問} & \textbf{答え} \\
    \hline
    第1問 & (1) & $x=\pm\sqrt{1-\sin y}$；グラフは解答の図を参照 \\
    \cline{2-3}
    & (2) & $4(\sqrt{2}-1)$ \\
    \hline
    第2問 & (1) & $h_0=\dfrac{\sqrt{6}}{3}$ \\
    \cline{2-3}
    & (2) & $t=\dfrac{\sqrt{3}}{4}\left(-h^3+\dfrac{\sqrt{6}}{2}h^2+h\right)$ \\
    \hline
    第3問 & -- & 左の等号は3ベクトルが互いに直交するとき，右の等号は同一平面上にあるとき \\
    \hline
    第4問 & (1)，(2) &  \\
    \hline
    第5問 & (1) & $E(X_1+Y_1)=n+1$ \\
    \cline{2-3}
    & (2) & $E(X_1)=\dfrac{2(n+1)}{3}$ \\
    \cline{2-3}
    & (3) & $E(Y_2)=\dfrac{n+1}{3}$ \\
    \hline
    第6問 & (1)，(2) & 証明問題 \\
    \hline
  \end{tabular}
\end{table}

\paragraph{試験形態}
6問．試験時間は要確認．

\end{document}
